\begin{center}
{\Large \bfseries \textcolor{lyred}{第七章 微分方程}}
\end{center}

\vspace{6pt}
{\large \bfseries \textcolor{lyred}{第7.1 节 微分方程的基本概念}}

\vspace{4pt}
{\bfseries \textcolor{lyred}{一.微分方程的例子}}

\vspace{4pt}
{\bfseries \textcolor{lyred}{二.微分方程的概念}}

{\bfseries \textcolor{lyred}{1.微分方程}}

表示未知函数,未知函数的导数及自变量之间关系的等式称为\textcolor{lyred}{微分方程}; 
微分方程中出现的未知函数导数的最高阶数称为该方程的\textcolor{lyred}{阶}. 
\textcolor{lyred}{例}.
y
x
'=
(一阶),
0.4
s''=-
(二阶),
sin
x y
xy
xy
y
x
'''
''
'
+
-
=
(三阶). 
\textcolor{lyred}{一阶方程}:
(
)
, ,
F x y y'=
,或者
(
)
,
y
f x y
'=
. 
\textcolor{lyred}{n 阶方程}:
()
(
)
, ,
,
n
F x y y
y
'
=
\Lambda ,
,或者
()
(
)
(
)
, ,
,
,
n
n
y
f x y y
y
-
'
=
\Lambda 
. 
{\bfseries \textcolor{lyred}{2.通解与特解}}

\textcolor{lyred}{通解}:含有n 个相互独立的任意常数项的解. 
\textcolor{lyred}{特解}:不包含任意常数项的解. 
\textcolor{lyred}{例}.
y
x
C
=
+
是
y
x
'=
的通解,
y
x
=
+是它的特解. 
\textcolor{lyred}{例}.
0.2
s
t
C t
C
=-
+
+
是
0.4
s''=-
的通解,
0.2
s
t
t
=-
+
是它的特解. 
\textcolor{lyred}{积分曲线}:方程的特解在坐标平面上代表的曲线. 
\textcolor{lyred}{例}.求方程
y
x
''=
的一条积分曲线,使它与
y
x
=
在(
)
2,4 处相切. 
解.
y
x
y
xdx
x
C
''
'
=
\Rightarrow 
=
=
+
\int 
,
(
)
y
x
C
dx
x
C x
C
=
+
=
+
+
\int 
,代入 
()
y
=
,
()
y'
=
,得
C =-,
C =
,故
y
x
x
=
-
+
. 
{\bfseries \textcolor{lyred}{3.初值问题}}

\textcolor{lyred}{例}.
()
y
x
y
'=


=

,其中()1
y
=
称为\textcolor{lyred}{初始条件}. 
\textcolor{lyred}{例}.
()
()
0.4
s
s
s
''=-


=

'
=

,其中()
s
=
, ()
s'
=
称为\textcolor{lyred}{初始条件}. 
\textcolor{lyred}{一阶初值问题}:
(
)
,
x x
y
f x y
y
y
=
'

=

=

,\textcolor{lyred}{二阶初值问题}:
(
)
, ,
x x
x x
y
f x y y
y
y
y
y
=
=
''
'
=

=

'
=

. 
类似地,还有\textcolor{lyred}{n 阶初值问题}. 
\vspace{6pt}
{\large \bfseries \textcolor{lyred}{第7.2 节 可分离变量的微分方程}}

\textcolor{lyred}{定义}.形如()
()
g y dy
f x dx
=
的方程称为\textcolor{lyred}{可分离变量}的微分方程. 
\textcolor{lyred}{解法}:由
()
()
g y dy
f x dx
=
\int 
\int 
,解得
()
()
G y
F x
C
=
+
,称为\textcolor{lyred}{隐式通解}. 
\textcolor{lyred}{例}.求
dy
xy
dx =
的通解. 
解.
ln
x
dy
xdx
y
x
C
y
Ce
y =
\Rightarrow 
=
+
\Rightarrow 
=
\int 
\int 
,其中
C
C
e
=\pm 
或0 . 
\textcolor{lyred}{注}.一般地,
()
()
ln
f x
y
f x
C
y
Ce
=
+
\Rightarrow 
=
. 
\textcolor{lyred}{例}.(人口模型)考虑人口函数
()
N t ,忽略政治,经济,环境等因素,人口增长率与 
人口基数成正比,即
ln
dN
dN
rN
rdt
N
rt
C
dt
N
=
\Rightarrow 
=
\Rightarrow 
=
+
\int 
\int 
,即
rt
N
Ce
=
,代入 
t =
,
N
N
=
,得
C
N
=
,故
r t
N
N e \cdots 
=
. 
\textcolor{lyred}{例}.(阻滞增长人口模型)设地球资源能容纳的最大人口数为K ,则 
(
)
(
)
1 ln
dN
dN
N
r K
N
rdt
rt
C
N dt
N K
N
K
K
N
=
-
\Rightarrow 
=
\Rightarrow 
=
+
-
-
\int 
\int 
,即 
rK t
N
Ce
K
N
\cdots 
=
-
,由
t =
,
N
N
=
,得
N
C
K
N
=
-
,故
rK t
N
N
e
K
N
K
N
\cdots 
=
-
-
,即 
(
)
rK t
KN
N
N
K
N
e-
\cdots 
=
+
-
,称为\textcolor{lyred}{Logistic 函数}. 
\textcolor{lyred}{例}.质量为m 的降落伞在空中由静止开始下落,设下落过程中的空气阻力与速度 
成正比,试求下落速度与下落时间的关系. 
解.设下落t 秒后速度
()
v
v t
=
,则
dv
dv
m
mg
kv
dt
dt
mg
kv
m
=
-
\Rightarrow 
=
-
\int 
\int 
,得 
ln mg
kv
t
C
k
m
-
-
=
+
,即
k t
m
mg
kv
Ce
-
\cdots 
-
=
,代入()
v
=
,得C
mg
=
,故 
k t
m
mg
v
e
k
-
\cdots 


=
-




,当t \to +\infty 时,
mg
v
k
\approx 
. 
\textcolor{lyred}{补充练习} 
1.求
y
xy
y
x
'+
-
=-
的通解. 
解.
(
)(
)
(
)
arctan
dy
y
y
x
x dx
y
x
x
C
y
'=
+
-
\Rightarrow 
=
-
\Rightarrow 
=
-
+
+
. 
2.求连续函数
()
f x ,使它满足
()
(
)
x
x
f t dt
x
tf x
t dt
=
+
-
\int 
\int 
. 
解.令u
x
t
=
-,则
(
)
(
)
()
()
()
x
x
x
x
tf x
t dt
x
u f u du
x
f u du
uf u du
-
=
-
=
-
\int 
\int 
\int 
\int 
, 
故
()
()
()
()
()
x
x
x
x
f t dt
x
x
f u du
uf u du
f x
f u du
=
+
-
\Rightarrow 
=+
\Rightarrow 
\int 
\int 
\int 
\int 
()
()
f
x
f x
'
=
,即
()
y
f x
=
满足dy
y
dx =
,解得
x
y
Ce
=
,由
x =
,
y =,得
C =,故 
()
x
f x
e
=
. 
\vspace{6pt}
{\large \bfseries \textcolor{lyred}{第7.3 节 齐次方程}}

\vspace{4pt}
{\bfseries \textcolor{lyred}{一.变量替换}}

\textcolor{lyred}{例}.求
dy
dx
x
y
=
+
+的通解. 
解.引入未知函数
u
x
y
=
+
+,
du
dy
dx
dx
=+
,代入方程,得
du
dx
u
=+
,解得 
ln
u
u
x
C
-
+
=
+
,即
ln
x
y
y
C
+
+
=
+
,故
y
x
y
Ce
+
+
=
. 
\textcolor{lyred}{例}.求
(
)
ln
dy
x
y
xy
dx -
-
=




的通解. 
解.
(
)
(
)
ln
d xy
y
xy
dx
-
=
,令u
xy
=
,
ln
du
u
u
dx
x
=
,解得
ln ln
ln
u
x
C
=
+
,即 
ln
Cx
u
Cx
u
e
=
\Rightarrow 
=
,故
Cx
xy
e
=
. 
\vspace{4pt}
{\bfseries \textcolor{lyred}{二.齐次方程}}

\textcolor{lyred}{定义}.形如dy
y
dx
x
\varphi 

=




的方程称为\textcolor{lyred}{齐次方程}. 
\textcolor{lyred}{例}.(
)
(
)
y
y
x
x
dy
x y
xy
x y
xy
dx
x
y
dy
dx
x
y
y
x


-

-


-
-
+
=
\Leftrightarrow 
=
=
+


+



. 
\textcolor{lyred}{解法}:令
y
u
x
=
,即y
xu
=
, dy
du
u
x
dx
dx
=
+
,代入方程,得
()
du
u
x
u
dx
\varphi 
+
=
,分离变量. 
\textcolor{lyred}{例}.求
2 dy
dy
y
x
xy
dx
dx
+
=
的通解. 
解.
y
dy
y
x
y
dx
xy
x
x






=
=
-
-
,令
y
u
x
=
,则y
xu
=
, dy
du
u
x
dx
dx
=
+
,代入方程,得 
du
u
u
x dx
u
+
=
-,解得
ln
ln
u
u
x
C
-
=
+
,即
ln
y
y
C
x
=
+
,故
y
x
y
Ce
=
. 
\textcolor{lyred}{例}.求
y
x
x
y
y =
+
+
'
的通解. 
解.以y 为自变量,得 
x
x
y
dx
x
x
dy
y
y
y
+
+
=
=
+
+,令
x
u
y
=
,则x
yu
=
, dx
du
u
y
dy
dy
=
+
,代入方程,得 
(
)
ln
ln
du
du
dy
u
y
u
u
u
u
y
C
dy
y
u
+
=
+
+\Rightarrow 
=
\Rightarrow 
+
+
=
+
+
\int 
\int 
,得 
(
)
u
u
Cy
u
Cy
u
+
+=
\Rightarrow 
+=
-
,故
x
y
C
C
=
+
. 
\textcolor{lyred}{补充练习 }
1.求
xy
y
xy
'+
=
的通解. 
解.
y
y
y
x
x
'+
=
,令
y
u
x
=
, y
xu
=
,则dy
du
u
x
dx
dx
=
+
,代入方程,得 
(
)
du
du
dx
u
x
u
u
dx
x
u
u
+
+
=
\Rightarrow 
=
-
,解得
xy
x
C
=
+
. 
\vspace{6pt}
{\large \bfseries \textcolor{lyred}{第7.4 节 一阶线性微分方程}}

\vspace{4pt}
{\bfseries \textcolor{lyred}{一.一阶线性方程}}

\textcolor{lyred}{一阶线性方程}:
()
()
dy
P x y
Q x
dx +
=
,其中
()
P x 为系数,
()
Q x 称为非齐次项; 
若
()
Q x \equiv 
,则称为\textcolor{lyred}{齐次的线性方程};否则称为\textcolor{lyred}{非齐次的线性方程}. 
\textcolor{lyred}{常数变易法}.(一)求对应的齐次线性方程
()
dy
P x y
dx +
=
的通解: 
()
()
P x dx
dy
P x dx
y
Ce
y
-\int 
=-
\Rightarrow 
=
,其中
()
P x dx
\int 
表示
()
P x 的某个原函数; 
(二)令
()
()
P x dx
y
u x e
-\int 
=
,代入方程,得 
()
()
P x dx
du e
Q x
dx
-\int 
=
,故有\textcolor{lyred}{通解公式}:
()
()
()
P x dx
P x dx
y
e
Q x e
dx
C
-


\int 
\int 
=
+




\int 
. 
\textcolor{lyred}{例}.求
(
)
dy
y
x
dx
x
-
=
+
+
的通解. 
解.(一)
(
)
dy
y
dy
dx
y
C x
dx
x
y
x
-
=
\Rightarrow 
=
\Rightarrow 
=
+
+
+
; 
(二)令
(
)
y
u x
=
+
,代入方程,得
(
)
(
)
du
x
x
dx \cdots 
+
=
+
,解得
(
)
u
x
C
=
+
+
,故 
(
)
(
)
y
x
C x
=
+
+
+
. 
\textcolor{lyred}{例}.求
cos
sin 2
y
x
y
y
'=
+
的通解. 
解.以y 为自变量,则
cos
sin 2
dx
y x
y
dy -
\cdots 
=
,故 
(
)
(
)
cos
cos
sin
sin 2
2sin
y dy
y dy
y
x
e
y e
dy
C
y
Ce
-
-
-


\int 
\int 
=
\cdots 
+
=-
-
+




\int 
. 
\textcolor{lyred}{例}.容器内盛有盐水100L ,初始含盐量50g ,现注入浓度
2g L
c =
的盐水,速度为 
3L min
v =
,同时又以
2L min
v =
的速度放水,求容器内的含盐量的变化规律. 
解.设t 分钟后含盐量为x 克,盐水浓度
x
c
t
=
+,故单位时间含盐量的变化为 
2 3
dx
x
dt
t
=
\cdots -
\cdots 
+
(流入-流出),即
dx
x
dt
t
+
=
+
,由通解公式,得 
(
)
(
)
2 100
dt
dt
t
t
C
x
e
e
dt
C
t
t
-
+
+


\int 
\int 
=
\cdots 
+
=
+
+


+


\int 
,其中
1.5 10
C =-
\times 
. 
\vspace{4pt}
{\bfseries \textcolor{lyred}{二.伯努利方程}}

\textcolor{lyred}{例}.求
ln
dy
y
y
x
dx
x
+
=
的通解. 
解.
(
)
2ln
2ln
d y
dy
y
y
y
x
x
dx
x
dx
x
-
-
-
-
+
=
\Leftrightarrow -
+
=
,令
z
y-
=
,得线性方程 
2ln
dz
z
x
dx
x
-
\cdots 
=-
,解得
(
)
ln
z
x C
x
=
-
,故
(
)
ln
y
x C
x
=
-
. 
\textcolor{lyred}{定义}.形如
()
()
(
)
0,1
n
dy
P x y
Q x y
n
dx +
=
\neq 
的方程称为\textcolor{lyred}{伯努利方程}. 
\textcolor{lyred}{解法}:
()
()
()
()
n
n
n
n
dy
d
y
P x y
Q x
y
P x y
Q x
dx
n dx
-
-
-
-
+
=
\Leftrightarrow 
\cdots 
+
=
-
,令
1 n
z
y -
=
,则 
(
)()
(
)
()
dz
n P x z
n Q x
dx +
-
=
-
,为线性方程,可解出z . 
\textcolor{lyred}{例}.设
()
f x 连续,且满足
()
()
()
x
f t
f x
dt
t f t
t
=+
+
\int 
,求
()
f x . 
解.
()
()
()
f x
f
x
x f x
x
'
=
+
,即
()
y
f x
=
满足
dy
y
dx
x y
x
=
+
,以y 为自变量,得 
dx
x
x
dy
y
=
+
,即
dx
dx
x
x
x
dy
y
dy
y
-
-
-
=
\Rightarrow 
+
=-,解得1
y
C
x
y
=-
+
; 
代入
x =,
y =,得
C =
,故
()
1 1
y
x
f x
x
y
x
+-
=-
+
\Rightarrow 
=
. 
\textcolor{lyred}{补充练习 }
1.设曲线
()
y
y x
=
过(
)
1,1 ,它在(
)
,x y 处的切线与坐标轴及过切点平行于y 轴的 
直线所围成的梯形面积等于常数
3a ,求
()
y
y x
=
. 
解.
y x
y
y x
a
'
-
\cdots +
+
\cdots 
=
,即
dy
a
y
dx
x
x
-
=-
,解得
2a
y
Cx
x
=
+
; 
代入
x =,
y =,得
C
a
=-
,故
(
)
1 2
a
y
a
x
x
=
+
-
. 
2.设
()
f x 连续,且满足
()
(
)
sin
x
f x
x
f x
t dt
=
-
-
\int 
,求
()
f x . 
解.
()
()
sin
x
f x
x
f u du
=
-\int 
,求导得
()
()
cos
f
x
x
f x
'
=
-
,即
()
y
f x
=
满足 
cos
y
y
x
'+
=
,解得
()
(
)
1 cos
sin
x
f x
x
x
Ce-
=
+
+
,
()
f
C
=
\Rightarrow 
=-
. 
3.设
()
f x 可导,且
()
()
(
)
f tx dt
f x
x
=
+
>
\int 
,
()1
f
=,求
()
f x . 
解.
()
()
x
f u du
xf x
x
=
+
\int 
,求导得
()
()
()
f x
f x
xf
x
'
=
+
+,即 
()
()
f
x
f x
x
x
'
+
=-
,解得
()
C
f x
x
=
-,
()1
f
C
=\Rightarrow 
=
. 
4.设
()
f x 连续,且满足
()
()
x
x
x
f x
e
e
f
t dt
=
-\int 
,求
()
f x . 
解.
()
()
()
()
()
x
x
x
x
x
f
x
e
e
f
t dt
e f
x
f x
e f
x
'
=
-
-
=
-
\int 
,即
()
y
f x
=
满足 
x
y
y
e y
'-
=-
,解得
()
x
x
f x
e
Ce-
=
+
,
()
f
C
=\Rightarrow 
=. 
\vspace{6pt}
{\large \bfseries \textcolor{lyred}{第7.5 节 可降阶的高阶微分方程}}

\textcolor{lyred}{类型一}.
()
()
n
y
f x
=
. 
\textcolor{lyred}{解法}:
(
)
()
()
n
n
y
y
dx
f
x dx
-=
=
\int 
\int 
,连续积分n 次.\textcolor{lyred}{ }
\textcolor{lyred}{例}.求
cos
x
y
e
x
'''=
-
的通解. 
解.积分得
sin
x
y
e
x
C
''=
-
+
,再积分得
cos
x
y
e
x
C x
C
'=
+
+
+
,再积分得 
sin
x
y
e
x
C x
C x
C
=
+
+
+
+
. 
\textcolor{lyred}{类型二}.
(
)
,
y
f x y
''
'
=
,缺y .\textcolor{lyred}{ }
\textcolor{lyred}{解法}:设dy
p
dx =
,则dp
y
dx
''
=
,代入方程,得
(
)
,
dp
f x p
dx =
,解出通解
(
)
,
p
p x C
=
,即 
(
)
,
dy
p x C
dx =
,于是
(
)
,
y
p x C
dx
C
=
+
\int 
. 
\textcolor{lyred}{例}.求(
)
x
y
xy
''
'
+
=
的通解. 
解.设dy
p
dx =
,则(
)
(
)
dp
dp
x
x
xp
dx
p
C
x
dx
p
x
+
=
\Rightarrow 
=
\Rightarrow 
=
+
+
,即 
(
)
1 1
dy
C
x
dx =
+
,故
(
)
(
)
y
C
x
dx
C
x
x
C
=
+
=
+
+
\int 
. 
\textcolor{lyred}{类型三}.
(
)
,
y
f
y y
''
'
=
,缺x .\textcolor{lyred}{ }
\textcolor{lyred}{解法}:设
dy
p
dx
=
,以y 为自变量,则
dp
dp dy
dp
y
p
dx
dy dx
dy
''=
=
\cdots 
=
,代入方程,得 
(
)
,
dp
p
f
y p
dy =
,解出
(
)
,
p
p y C
=
,即
(
)
,
dy
p y C
dx =
,再分离变量,得 
(
)
,
dy
dx
p y C
=
,解出
(
)
,
dy
x
C
p y C
=
+
\int 
. 
\textcolor{lyred}{例}.求
yy
y
''
'
-
=
的通解. 
解.设
dy
p
dx
=
,以y 为自变量,则
dp
dp dy
dp
y
p
dx
dy dx
dy
''=
=
\cdots 
=
,代入方程,得 
ln
ln
dp
dp
dy
yp
p
p
y
C
p
C y
dy
p
y
-
=
\Rightarrow 
=
\Rightarrow 
=
+
\Rightarrow 
=
,即
dy
C y
dx =
\Rightarrow  
ln
C x
dy
C dx
y
C x
C
y
C e
y
=
\Rightarrow 
=
+
\Rightarrow 
=
\int 
\int 
. 
\textcolor{lyred}{例}.设地球质量为M ,半径为R ,一质量为m 的火箭,由地面以速度
2GM
v
R
=
垂直向上发射,求火箭高度r 与时间t 的关系. 
解.
(
)
d r
GMm
m dt
R
r
=-
+
,令
dr
v
dt
=
,以r 为自变量,则
d r
dv dr
dv
v
dt
dr dt
dr
=
\cdots 
=
, 
代入方程,得
(
)
(
)
dv
GM
GM
v
vdv
dr
dr
R
r
R
r
=-
\Rightarrow 
=
+
+
,解得
v
GM
C
R
r
=
+
+
, 
代入
r =
,
2GM
v
R
=
,得
C =
,于是
2GM
v
R
r
=
+
,即
2GM
v
R
r
=
+
,故 
dr
GM
R
rdr
GM dt
dt
R
r
=
\Rightarrow 
+
=
+
,解得(
)
3 R
r
GM t
C
+
=
\cdots +
, 
代入
t =
,
r =
,得
C
R
=
,因此(
)
R
r
GM t
R
+
=
\cdots +
. 
\vspace{6pt}
{\large \bfseries \textcolor{lyred}{第7.6 节 高阶线性微分方程}}

\textcolor{lyred}{二阶齐次线性方程}:
()
()
y
P x y
Q x y
''
'
+
+
=
-----(1); 
\textcolor{lyred}{二阶非齐次线性方程}:
()
()
()
y
P x y
Q x y
f x
''
'
+
+
=
-----(2). 
\textcolor{lyred}{定理}.设
()
1y
x 与
()
2y
x 均是(1)的解,则
()
()
y
C y
x
C y
x
=
+
也是(1)的解,其中 
C ,
C 为任意常数. 
\textcolor{lyred}{定义}.若存在不全为零的常数
,
,
,
n
k k
k
\Lambda 
,使
n
n
k y
k y
k y
+
+
+
\equiv 
\Lambda 
,则称这n 个 
函数\textcolor{lyred}{线性相关},否则称为\textcolor{lyred}{线性无关}. 
\textcolor{lyred}{定理}.设
()
1y
x 与
()
2y
x 是(1)的两个线性无关的特解,则
()
()
y
C y
x
C y
x
=
+
是 
(1)的通解. 
\textcolor{lyred}{例}.已知
y
y
y
''
'
-
+
=
有两个线性无关的解
x
y
e
=
和
x
y
xe
=
,故它的通解为 
x
x
y
C e
C xe
=
+
. 
\textcolor{lyred}{注}.一般地,若函数
()
()
()
,
,
,
n
y
x
y
x
y
x
\Lambda 
是n 阶齐次线性微分方程 
()
()
(
)
()
()
n
n
n
n
y
a
x y
a
x y
a
x y
-
-
'
+
+
+
+
=
\Lambda 
的n 个线性无关的特解,则 
()
()
()
n
n
y
C y
x
C y
x
C y
x
=
+
+
+
\Lambda 
是该方程的通解. 
\textcolor{lyred}{定理}.设()
y x 是(1)的解,
()
y
x
*
是(2)的解,则
()
()
y
y x
y
x
*
=
+
是(2)的解. 
\textcolor{lyred}{定理}.设(
)
;
,
y x C C
是(1)的通解,
()
y
x
*
是(2)的特解,则y
y
y*
=
+
是(2)的通解. 
\textcolor{lyred}{例}.
y
y
y
''
'
-
+
=
有特解
y*=,故
x
x
y
C e
C xe
=
+
+为它的通解. 
\textcolor{lyred}{定理}.设
()
ky
x
*
是
()
()
()(
)
1,2
k
y
P x y
Q x y
f
x
k
''
'
+
+
=
=
的解,则 
()
()
y
y
x
y
x
*
*
=
\pm 
是
()
()
()
()
y
P x y
Q x y
f
x
f
x
''
'
+
+
=
\pm 
的解. 
\textcolor{lyred}{推论}.设
()
1y
x
*
与
()
2y
x
*
均是(2)的解,则
()
()
y
x
y
x
*
*
-
是(1)的解. 
\textcolor{lyred}{例}.设
()
1y
x ,
()
2y
x 及
()
3y
x 均为(2)的解,则该方程必有解______. 
(A)
y
y
y
+
+
   (B)
y
y
y
+
-
   (C)
y
y
y
-
-
   (D)
y
y
y
-
-
-
解.
(
)
y
y
y
y
y
y
+
-
=
-
+
是(2)的解,故选(B). 
\textcolor{lyred}{例}.已知
y =
,
y
x
=
+
,
x
y
x
e
=
+
+
均为(2)的解,求通解. 
解.
(
)
(
)
x
y
C
y
y
C
y
y
y
C e
C x
=
-
+
-
+
=
+
+. 
\vspace{6pt}
{\large \bfseries \textcolor{lyred}{第7.7 节 常系数齐次线性微分方程}}

\vspace{4pt}
{\bfseries \textcolor{lyred}{一.二阶常系数齐次线性方程}}

(1)
y
p y
q y
''
'
+
+
=
,其中p ,q 为常数. 
\textcolor{lyred}{特征方程法}:由于
r x
e \cdots 和它的各阶导数之间只差一个常数倍,故假设(1)具有形如 
r x
y
e \cdots 
=
的特解,其中r 待定,则
(
)
rx
y
p y
q y
r
pr
q e
''
'
+
+
=
\Leftrightarrow 
+
+
=
\Leftrightarrow  
r
pr
q
+
+
=
,解出
1,2
p
p
q
r
-
\pm 
-
=
,得(1)的两个特解
r x
y
e
=
,
r x
y
e
=
. 
\textcolor{lyred}{定义}.(2)
r
pr
q
+
+
=
,称为(1)的\textcolor{lyred}{特征方程},
1,2
p
p
q
r
-
\pm 
-
=
称为\textcolor{lyred}{特征根}. 
\textcolor{lyred}{情形1}.
\Delta >
,即(2)有两个不相等的实根1r 与2r ,此时(1)有两个线性无关的解 
r x
y
e
=
,
r x
y
e
=
,通解
r x
r x
y
C e
C e
=
+
; 
\textcolor{lyred}{情形2}.
\Delta =
,即(2)有两个相等实根1
p
r
r
=
=-
,此时(1)有一个解
r x
y
e
=
,而 
r x
y
xe
=
是另一个解,通解
(
)
r x
y
C
C x e
=
+
; 
\textcolor{lyred}{情形3}.
\Delta <
,即(2)有一对共轭复根1,2
r
i
\alpha 
\beta 
=
\pm 
,此时(1)有两个线性无关的解 
cos
x
e
x
\alpha 
\beta ,
sin
x
e
x
\alpha 
\beta ,通解
(
)
cos
sin
x
y
e
C
x
C
x
\alpha 
\beta 
\beta 
=
+
. 
\textcolor{lyred}{例}.求
y
y
y
''
'
-
-
=
的通解. 
解.
(
)(
)
r
r
r
r
r
-
-
=
\Rightarrow 
-
+
=
\Rightarrow 
=-,3,故
x
x
y
C e
C e
-
=
+
. 
\textcolor{lyred}{例}.求
y
y
y
''
'
+
+
=
的通解. 
解.
(
)
r
r
r
r
+
+=
\Rightarrow 
+
=
\Rightarrow 
=-, 1
-,故
(
)
x
y
C
C x e-
=
+
. 
\textcolor{lyred}{例}.求
y
y
y
''
'
-
+
=
的通解. 
解.
(
)
r
r
r
r
i
-
+
=
\Rightarrow 
-
=-\Rightarrow 
=\pm 
,故
(
)
cos2
sin 2
x
y
e
C
x
C
x
=
+
. 
\textcolor{lyred}{例}.求以
(
)
x
y
C
C x e
=
+
为通解的二阶齐次线性微分方程. 
解.特征方程为(
)
r
r
r
-
=
-
+
,故所求方程为
y
y
y
''
'
-
+
=
. 
\textcolor{lyred}{例}.求
y
y
''+
=
的在(
)
, 1
\pi -
处和直线
y
x
\pi 
+=
-
相切的积分曲线. 
解.
r
r
i
+
=
\Rightarrow 
=\pm \cdots ,通解
cos3
sin3
y
C
x
C
x
=
+
,由
x
x
y
y
\pi 
\pi 
=
=

=-
'
=

,得 
cos3
sin3
1,
sin3
cos3
C
C
C
C
C
C
\pi 
\pi 
\pi 
\pi 
+
=-

\Rightarrow 
=
=-
-
+
=

,故
cos3
sin3
y
x
x
=
-
. 
\vspace{4pt}
{\bfseries \textcolor{lyred}{二.n 阶常系数齐次线性方程}}

(3)
()
(
)1
n
n
n
n
y
p y
p
y
p y
-
-
'
+
+
+
+
=
\Lambda 
,其中
,
,
,
n
p
p
p
\Lambda 
为常数. 
(4)
n
n
n
n
r
p r
p
r
p
-
-
+
+
+
+
=
\Lambda 
称为(3)的\textcolor{lyred}{特征方程}. 
若r 是(4)的k 重实根,则(3)有一组特解
m
r x
y
x e
=
,0
m
k
\le 
\le 
-; 
若r
i
\alpha 
\beta 
=
\pm 
是(4)的k 重共轭复根,则(3)有两组特解: 
cos
m
x
y
x e
x
\alpha 
\beta 
=
,
sin
m
x
y
x e
x
\alpha 
\beta 
=
,0
m
k
\le 
\le 
-. 
\textcolor{lyred}{例}.求
y
y
y
y
'''
''
'
-
+
+
=
的通解. 
解.
(
)(
)
(
)(
)(
)
r
r
r
r
r
r
r
r
r
-
+
+
=
\Rightarrow 
-
+
+
=
\Rightarrow 
-
-
+
=
,得 
r =-, 2 ,5,故
x
x
x
y
C e
C e
C e
-
=
+
+
. 
\textcolor{lyred}{例}.求
()
()
()
y
y
y
y
y
y
''
'
+
+
+
+
+
=
的通解. 
解.
(
)(
)
r
r
r
r
r
r
r
+
+
+
++=
\Rightarrow 
+
+
=
,得
r =-, i
\pm , i
\pm ,故 
cos
sin
cos
sin
x
y
C e
C
x
C
x
C x
x
C x
x
-
=
+
+
+
+
. 
\textcolor{lyred}{例}.求具有特解
x
y
xe
=
及
x
y
e-
=
的三阶常系数齐次线性微分方程. 
解.(
)(
)
r
r
r
r
r
-
+
=
\Rightarrow 
-
-+=
,故方程为
y
y
y
y
'''
''
'
-
-
+
=
. 
\textcolor{lyred}{补充练习 }
1.求
()
y
y
y
'''
''
-
+
=
的通解. 
解.
(
)
r
r
r
r
r
r
r
-
+
=
\Rightarrow 
-
+
=
\Rightarrow 
=
,0 ,1
2i
\pm 
,故 
cos2
sin 2
x
x
y
C
C x
C e
x
C e
x
=
+
+
+
. 
2.利用变换
cos
x
t
=
化简方程(
)
x
y
xy
y
''
'
-
-
+
=
,并求通解. 
解.
sin
dy dt
dy
y
dt dx
t dt
'=
=-
,
cos
sin
sin
sin
dy dt
t dy
d y
y
dt dx
t dt
t dt
t
'



''=
=
-
-







, 
代入方程,得
d y
y
dt
+
=
,解得
cos
sin
y
C
t
C
t
C x
C
x
=
+
=
+
-
. 
3.设
()
y
y x
=
满足方程
y
by
cy
''
'
+
+
=
,其中,
b c >
,求
()
lim
x
y x
\to +\infty 
. 
解.
r
br
c
+
+
=
,
1,2
b
b
c
r
-\pm 
-
=
. 
(1)当
b
c
-
>
时,
r x
r x
y
C e
C e
=
+
,而1,2
r
<
,故
()
lim
x
y x
\to +\infty 
=
; 
(2)当
b
c
-
=
时,
(
)
b x
y
C
C x e
-
=
+
,而
r <
,故
()
lim
x
y x
\to +\infty 
=
; 
(3)当
b
c
-
<
时,
(
)
cos
sin
b x
y
e
C
x
C
x
\beta 
\beta 
-
=
+
,故
()
lim
x
y x
\to +\infty 
=
. 
\vspace{6pt}
{\large \bfseries \textcolor{lyred}{第7.8 节 常系数非齐次线性方程}}

(1)
()
y
py
qy
f x
''
'
+
+
=
,其中p ,q 为常数. 
由于叠加原理,求(1)的通解等价于求对应的齐次线性方程的通解(
)
;
,
y x C C
和 
(1)的一个特解
()
y
x
*
,而求y*的方法为\textcolor{lyred}{待定系数法}. 
\textcolor{lyred}{情形一}.
()
()
x
m
f x
P
x e\lambda 
=
,其中
()
m
P
x 为m 次多项式. 
设
()
k
x
m
y
x Q
x e\lambda 
*=
,其中
()
m
Q
x 为待定的m 次多项式,当\lambda 不是特征根时,
k =
, 
当\lambda 是单特征根时,
k =,当\lambda 是重特征根时,
k =
. 
\textcolor{lyred}{例}.求
(
)
x
y
y
y
x
e
''
'
+
-
=
+
的通解. 
解.(一)
r
r
r
+
-
=
\Rightarrow 
=-, 1
2 ,故
x
x
y
C e
C e
-
=
+
; 
(二)
\lambda =不是特征根,令
(
)
*
x
y
ax
bx
c e
=
+
+
,代入原方程,得 
(
)
ax
a
b x
a
b
c
x


+
+
+
+
+
=
+




,于是
a =,
b =-,
c =
,即 
*
x
y
x
x
e


=
-
+




,故
x
x
x
y
C e
C e
x
x
e
-


=
+
+
-
+




. 
\textcolor{lyred}{例}.求
(
)
x
y
y
y
x e-
''
'
-
-
=
-
的通解. 
解.(一)
r
r
r
-
-
=
\Rightarrow 
=-,2 ,故
x
x
y
C e
C e
-
=
+
; 
(二)
\lambda =-为单特征根,令
(
)
*
x
y
x ax
b e-
=
+
,代入原方程,得 
(
)
ax
a
b
x
-
+
-
=-
+,于是
a =,
b =-,即
(
)
*
x
y
x
x e-
=
-
,故 
(
)
x
x
x
y
C e
C e
x
x e
-
-
=
+
+
-
. 
\textcolor{lyred}{例}.求
(
)
x
y
y
y
x
e
''
'
-
+
=
+
的通解. 
解.(一)
r
r
r
-
+
=
\Rightarrow 
=
,3,故
(
)
x
y
C
C x e
=
+
; 
(二)
\lambda =
为重特征根,令
(
)
*
3x
y
x
ax
b e
=
+
,代入原方程,得6
ax
b
x
+
=
+,于是 
a =
,
b =
,即
*
x
y
x
x
e


=
+




,故
(
)
x
x
y
C
C x e
x
x
e


=
+
+
+




. 
\textcolor{lyred}{情形二}.
()
()
()
(
)
cos
sin
x
l
n
f x
e
P x
x
P
x
x
\lambda 
\omega 
\omega 
\omega 
=
+
\neq 




,其中
()
lP x 为l 次多项式, 
()
nP x 为n 次多项式. 
设
()
()
cos
sin
k
x
m
m
y
x e
Q
x
x
R
x
x
\lambda 
\omega 
\omega 
*=
+



,
(
)
max
,
m
l n
=
,其中
()
m
Q
x ,
()
m
R
x 为 
待定的m 次多项式,当
i
\lambda 
\omega 
+
不是特征根时,
k =
;当
i
\lambda 
\omega 
+
是特征根时
k =. 
\textcolor{lyred}{例}.求
cos 2
y
y
x
x
''+
=
的通解. 
解.(一)
r
r
i
+=
\Rightarrow 
=\pm ,故
cos
sin
y
C
x
C
x
=
+
; 
(二)
i
i
\lambda 
\omega 
+
=
不是特征根,令
(
)
(
)
*
cos2
sin 2
y
ax
b
x
cx
d
x
=
+
+
+
,代入 
原方程,得(
)
(
)
cos2
sin 2
cos2
ax
b
c
x
cx
d
a
x
x
x
-
-
+
-
+
+
=
,于是 
a =-
,
b =
,
c =
,
d =
,即
*
cos2
sin 2
y
x
x
x
=-
+
,故 
cos
sin
cos2
sin 2
y
C
x
C
x
x
x
x
=
+
-
+
. 
\textcolor{lyred}{例}.求
cos2
x
y
y
y
e
x
''
'
-
+
=
的通解. 
解.(一)
r
r
r
i
-
+
=
\Rightarrow 
=\pm 
,故
(
)
cos2
sin 2
x
y
e
C
x
C
x
=
+
; 
(二)
i
i
\lambda 
\omega 
+
=+
是特征根,令
(
)
*
cos2
sin 2
x
y
x a
x
b
x e
=
+
,代入方程,得 
4 cos2
4 sin 2
cos2
b
x
a
x
x
-
=
,于是
a =
,
b =
,即
*
sin 2
x
y
xe
x
=
,故 
(
)
cos2
sin 2
sin 2
x
x
y
e
C
x
C
x
xe
x
=
+
+
. 
\textcolor{lyred}{例}.求
(
)
cos2
y
y
x
x
''+
=
+
的通解. 
解.(一)
r
r
i
+
=
\Rightarrow 
=\pm 
,故
cos2
sin 2
y
C
x
C
x
=
+
; 
(二)设
*
*
*
y
y
y
=
+
,其中(1)
*
1y
ax
b
=
+
,代入
y
y
x
''+
=
,得
*
y
x
=
; 
(2)
(
)
*
cos2
sin 2
y
x a
x
b
x
=
+
,代入
cos2
y
y
x
''+
=
,得
*
1 sin 2
y
x
x
=
,故 
1 sin 2
y
x
x
x
*=
+
,因此
cos2
sin 2
sin 2
y
C
x
C
x
x
x
x
=
+
+
+
. 
\textcolor{lyred}{补充练习 }
1.求以
(
)
x
x
y
C
x e
C e-
=
+
+
为通解的二阶常系数线性方程. 
解.
x
x
x
y
C e
C e
xe
-
=
+
+
,特征方程为(
)(
)
r
r
r
r
-
+
=
\Leftrightarrow 
+-
=
, 
故
()
y
y
y
f x
''
'
+
-
=
,代入
x
y
xe
*=
,得
()
3 x
f x
e
=
,故
x
y
y
y
e
''
'
+
-
=
. 
2.求以
(
)
cos
sin
x
y
C
x
C
x
e
=
+
+
为通解的二阶常系数线性方程. 
解.
(
)
cos
sin
x
x
y
e
C
x
C
x
e
=
+
+
,
r
i
=+是特征根,故特征方程为(
)
r -
=-,即 
r
r
-
+
=
,故
()
y
y
y
f x
''
'
-
+
=
,代入
x
y
e
*=
,得
()
x
f x
e
=
,故所求方程为 
x
y
y
y
e
''
'
-
+
=
. 
3.设
x
x
y
xe
e
=
+
,
x
x
y
xe
e-
=
+
,
x
x
x
y
xe
e
e-
=
+
+
均为一个二阶常系数非齐次 
线性方程的特解,求此方程. 
解.
x
y
y
e-
-
=
,
x
y
y
e
-
=
为对应齐次线性方程的解,故特征方程 
(
)(
)
r
r
-
+
=
,即
r
r
--
=
,故所求方程为
()
y
y
y
f x
''-
-
=
,又代入 
x
y
xe
=
,得
()
(
)
x
f x
e
x
=
-
. 
4.设
()
f x 有连续的二阶导数,
()
()
()
x
t
f x
f
t
f t
te
dt
-
''


=+
+
-


\int 
,且 
()
f '
=
,求
()
f x . 
解.
()
()
()
x
f
x
f
x
f x
xe-
'
''


=
+
-

,
()
y
f x
=
满足
x
y
y
y
xe-
''
'
-
+
=
; 
(一)
r
r
r
-
+
=
\Rightarrow 
=,2 ,故
x
x
y
C e
C e
=
+
; 
(二)
\lambda =-不是特征根,设
(
)
x
y
ax
b e
*
-
=
+
,代入原方程,得 
x
y
x
e
*
-


=
+




,故
x
x
x
y
C e
C e
x
e-


=
+
+
+




,由
x =
,
y =,
y'=
,得 
C =
,
C =-
,故
()
x
x
x
f x
e
e
x
e-


=
-
+
+




. 
5.求
sin
y
a y
x
''+
=
的通解,其中
a >
. 
解.
r
a
+
=
,r
ai
=\pm 
;(1)若
a \neq ,设
cos
sin
y
A
x
B
x
*=
+
,代入方程,得 
A =
,
B
a
=
-,故
cos
sin
sin
y
C
ax
C
ax
x
a
=
+
+
-
; 
(2)若
a =,设
(
)
cos
sin
y
x A
x
B
x
*=
+
,代入方程,得
A =-
,
B =
,故 
cos
sin
cos
y
C
x
C
x
x
x
=
+
-
. 
6.求
cos
y
y
x
''+
=
的通解. 
解.(一)
r
r
i
+
=
\Rightarrow 
=\pm 
,故
cos2
sin 2
y
C
x
C
x
=
+
; 
(二)
cos
cos2
x
x
=
+
,(1)
y
y
''+
=
,
y*=
;(2)
cos2
y
y
x
''+
=
, 
令
(
)
*
cos2
sin 2
y
x a
x
b
x
=
+
,代入方程,得
sin 2
x
y
x
*=
,故 
sin 2
x
y
x
*=
+
,于是
cos2
sin 2
sin 2
x
y
C
x
C
x
x
=
+
+
+
. 
7.求
(
)
sin
d x
dx
y
x
dy
dy


+
+
=




满足条件()
y
=
,
()
y'
=
的解. 
解.
dx
dy
y
=
',
d x
d
dx
y
y
dy
dx
y
dy
y
y
y
''
''


=
=-
=-

'
'
'
'


,代入方程,得 
sin
y
y
x
''-
=
,解得
1 sin
x
x
y
C e
C e
x
-
=
+
-
,由()
y
=
,
()
y'
=
,得 
C =,
C =-,故
1 sin
x
x
y
e
e
x
-
=
-
-
. 
8.利用变换
tan
x
u
=
化简方程(
)
(
)
d y
dy
x
x
x
x
y
dx
dx
x
+
+
+
+
=
+
,并求通解. 
解.
cos
dy
dy du
dy
u
dx
du dx
du
=
=
\cdots 
,
2cos sin
cos
cos
d y
dy
d y
u
u
u
u
dx
du
du


=-
\cdots 
+
\cdots 




, 
代入方程,得
sin
d y
y
u
du +
=
,解得
cos
sin
cos
y
C
u
C
u
u
u
=
+
-
,即 
arctan
2 1
C
C x
x
y
x
x
x
=
+
-
+
+
+
. 